% !TEX root = paper.tex
% Shared macros. Requires xcolor, amssymb, and bm (loaded in paper.tex).

% ---------------------------------------------------------------- Notation
\newcommand{\vect}[1]{\boldsymbol{#1}}
\newcommand{\R}{\mathbb{R}}
\newcommand{\method}{moving-patch MPM}

% ---------------------------------------------------------------- Drafting
% Red placeholder for results, numbers, and claims that still need measured evidence.
\newcommand{\pending}[1]{\textcolor{red}{[#1]}}
% Brown marker for provisional numbers.
\newcommand{\unchecked}[1]{\textcolor{brown}{#1}}

% ---------------------------------------------------------------- Figures
% Shows a figure once its file exists, and a TODO box until then. #1: extra \includegraphics options, #2: file.
% Requires graphicx (loaded in paper.tex).
\newcommand{\figureortodo}[2][]{%
  \IfFileExists{#2}%
    {\includegraphics[width=\linewidth,#1]{#2}}%
    {\fbox{\parbox[c][2.0cm][c]{0.95\linewidth}{\centering\junnosuke{TODO: figure \texttt{\detokenize{#2}}}}}}%
}

% ---------------------------------------------------------------- Tables
% Mean and standard deviation in a table cell: \ms{mean}{std}. The mean is black, and the "± std" part is grey and smaller.
% To drop the standard deviation everywhere, replace the definition by: \newcommand{\ms}[2]{#1}
\newcommand{\ms}[2]{#1{\tiny\,\textcolor{gray}{$\pm$\,#2}}}
% Table captions in normal upright type, like the figure captions. This departs from the IEEEtran default, which sets
% the table caption text in small capitals; the label ("TABLE I") is unchanged. To revert, delete the \patchcmd line.
\RequirePackage{etoolbox}
\makeatletter
\patchcmd{\@makecaption}{\scshape}{}{}{\PackageWarning{shortcut}{Could not patch the table caption font}}
\makeatother
% Placeholder for a table cell whose value is still to come.
\newcommand{\todocell}{\textcolor{blue}{\textbf{?}}}
% Arrows for table headers that show the better direction of a metric: higher is better, lower is better.
\newcommand{\higherbetter}{$\uparrow$}
\newcommand{\lowerbetter}{$\downarrow$}

% ---------------------------------------------------------------- Headings
% Sub-subsections: numbered B-1, B-2, C-1, ... (no parenthesis, no colon), with the title on its own line and without indent.
% The positive skip after the title makes it a display heading, so the text starts on a new line with the usual paragraph indent.
\renewcommand{\thesubsubsectiondis}{\Alph{subsection}-\arabic{subsubsection}}
\makeatletter
\def\subsubsection{\@startsection{subsubsection}{3}{\z@}{0.8ex plus 0.5ex minus 0.2ex}%
{0.3ex plus 0.2ex minus 0ex}{\normalfont\normalsize\itshape}}
\makeatother

% ---------------------------------------------------------------- Revision
\newcommand{\revision}[1]{\textcolor{teal}{\%#1}}
\newcommand*{\Revised}{\textcolor{black}}

% ---------------------------------------------------------------- Comments
% \max is a built-in math operator, so Max uses \maxcomment.
\newcommand{\junnosuke}[1]{\textbf{\textcolor{blue}{Jn: #1}}}
\newcommand{\junnosukereply}[1]{\textbf{\textcolor{blue}{[Jn: #1]}}}
\newcommand{\antoine}[1]{\textbf{\textcolor{purple}{Antoine: #1}}}
\newcommand{\antoinereply}[1]{\textbf{\textcolor{purple}{[Antoine: #1]}}}
\newcommand{\maxcomment}[1]{\textbf{\textcolor{orange}{Max: #1}}}
\newcommand{\maxreply}[1]{\textbf{\textcolor{orange}{[Max: #1]}}}
\newcommand{\mayank}[1]{\textbf{\textcolor{magenta}{Mayank: #1}}}
\newcommand{\mayankreply}[1]{\textbf{\textcolor{magenta}{[Mayank: #1]}}}
\newcommand{\kelly}[1]{\textbf{\textcolor{olive}{Kelly: #1}}}
\newcommand{\kellyreply}[1]{\textbf{\textcolor{olive}{[Kelly: #1]}}}
